
\subsection*{\SGARefTarget{sga7.sga7i.evii.n3-4}{3.4}. Variance of $\underline\Psi_{L_\bullet}$ with
$L_\bullet$}

Consider a morphism of complexes
\begin{equation}\tag{\SGARefTarget{sga7.sga7i.evii.n3-4-1}{3.4.1}}
 f:L'_\bullet\longrightarrow L_\bullet
\end{equation}
given by a commutative diagram
\begin{equation}\tag{\SGARefTarget{sga7.sga7i.evii.n3-4-2}{3.4.2}}
\begin{tikzcd}[column sep=large,row sep=large]
 {}\arrow[r]&L'_2\arrow[r]\arrow[d,"f_2"']
 &L'_1\arrow[r]\arrow[d,"f_1"']
 &L'_0\arrow[r]\arrow[d,"f_0"']&0\\
 {}\arrow[r]&L_2\arrow[r]&L_1\arrow[r]&L_0\arrow[r]&0.
\end{tikzcd}
\end{equation}
We associate with it a canonical functor
\begin{equation}\tag{\SGARefTarget{sga7.sga7i.evii.n3-4-3}{3.4.3}}
 f^*:\underline\Psi_{L_\bullet}\longrightarrow
 \underline\Psi_{L'_\bullet}.
\end{equation}
It is in fact a cocartesian functor of cofibred categories over
$\underline C$: it is compatible with the structural functors to
$\underline C$ and carries cocartesian morphisms to cocartesian
morphisms; see SGA 1, Exposé VI.

If $(X,\alpha)$ is an object of $\underline\Psi_{L_\bullet}$, define
$f^*(X,\alpha)=(X',\alpha')$ as follows. The extension $X'$ of
$L'_0$ by $A$ is the pullback, along $f_0$, of the given extension
$X$ of $L_0$ by $A$, and $\alpha'$ is the partial trivialization
deduced from $\alpha$ by transitivity in the second square of
(\SGARefLink{sga7.sga7i.evii.n3-4-2}{3.4.2}). The compatibility of $\alpha'$ with the relation
$d'_0d'_1=0$ follows immediately from the analogous compatibility
for $\alpha$. By construction, (\SGARefLink{sga7.sga7i.evii.n3-4-3}{3.4.3}) commutes with the structural
functors to $\underline C$. Its cocartesian character follows at
once from the fact that, in the category of extensions of objects of
$\underline C$, direct images of extensions commute with inverse
images of extensions.

The functor (\SGARefLink{sga7.sga7i.evii.n3-4-3}{3.4.3}) induces morphisms of functors
\begin{equation}\tag{\SGARefTarget{sga7.sga7i.evii.n3-4-4}{3.4.4}}
 f^i:\Psi^i_{L_\bullet}\longrightarrow\Psi^i_{L'_\bullet},
 \qquad i=0,1
\end{equation}
(compare \SGARefLink{sga7.sga7i.evii.n2-6}{2.6}). The morphism (\SGARefLink{sga7.sga7i.evii.n3-4-1}{3.4.1}), on the other hand, induces
morphisms
\begin{equation}\tag{\SGARefTarget{sga7.sga7i.evii.n3-4-5}{3.4.5}}
 f^i:\Ext^i(L_\bullet,-)\longrightarrow\Ext^i(L'_\bullet,-),
 \qquad i\in\mathbb Z.
\end{equation}
For $i=0,1$, the morphisms (\SGARefLink{sga7.sga7i.evii.n3-4-4}{3.4.4}) and (\SGARefLink{sga7.sga7i.evii.n3-4-5}{3.4.5}) are compatible with
the isomorphisms (\SGARefLink{sga7.sga7i.evii.n3-1-8}{3.1.8}) and (\SGARefLink{sga7.sga7i.evii.n3-2-4}{3.2.4}). In other words, for
$A\in\operatorname{Ob}\underline C$ and $i=0,1$, the following
diagram is commutative:
\begin{equation}\tag{\SGARefTarget{sga7.sga7i.evii.n3-4-6}{3.4.6}}
\begin{tikzcd}[column sep=large,row sep=large]
 \Psi^i_{L_\bullet}(A)\arrow[r,"\sim"]\arrow[d]
 &\Ext^i(L_\bullet,A)\arrow[d]\\
 \Psi^i_{L'_\bullet}(A)\arrow[r,"\sim"]
 &\Ext^i(L'_\bullet,A).
\end{tikzcd}
\end{equation}
The verification of this compatibility is essentially immediate and
is left to the reader.

\underline{Proposition} \SGARefTarget{sga7.sga7i.evii.n3-4-7}{3.4.7}. \underline{The functor} (\SGARefLink{sga7.sga7i.evii.n3-4-3}{3.4.3})
\underline{is faithful} (\underline{respectively fully faithful,
respectively an equivalence of categories}) \underline{if and only
if}
\[
 H_0(f):H_0(L'_\bullet)\longrightarrow H_0(L_\bullet)
\]
\underline{is an epimorphism} (\underline{respectively, if}
$H_0(f)$ \underline{is an isomorphism and} $H_1(f)$
\underline{is an epimorphism; respectively, if} $H_0(f)$
\underline{and} $H_1(f)$ \underline{are isomorphisms}).

Indeed, with $f^i$ denoting the functorial morphism (\SGARefLink{sga7.sga7i.evii.n3-4-4}{3.4.4}), the
required condition is that $f^0$ be a monomorphism (respectively,
that $f^0$ be an isomorphism and $f^1$ a monomorphism; respectively,
that $f^0$ and $f^1$ be isomorphisms). By the commutativity of
(\SGARefLink{sga7.sga7i.evii.n3-4-6}{3.4.6}), these conditions may be expressed by the analogous
conditions on the functors $\Ext^i$, $i=0,1$. The proposition then
reduces to a straightforward statement of homological algebra,
which we leave to the reader as an exercise.

